\section{Conclusion}

The ESS-corrected permutation test converts a redistricting ensemble percentile
into a court-admissible p-value. The key practical implications:

\begin{enumerate}
  \item \textbf{1{,}000-step ensembles are underpowered for North Carolina.}
    NC's ESS of 70 at 1K steps means the 0.003 percentile position corresponds
    to a borderline-significant $p = 0.041$.
    Courts should not treat 1{,}000-step NC ensemble results as conclusive.
  \item \textbf{10{,}000-step ensembles are reliable for most states.}
    NC, WI, GA, PA all achieve $p < 0.01$ at 10K steps.
    Texas requires 30K steps for similar reliability (S.3).
  \item \textbf{The \texttt{bisect-analysis::permutation} module provides
    a drop-in implementation.}
    Any user with a GerryChain ensemble JSON output and the G.4 ESS estimate
    can compute the calibrated p-value with three lines of Rust
    (or the equivalent Python/scipy Beta CDF computation from S.2).
  \item \textbf{Two-sided tests are appropriate for partisan fairness metrics.}
    The lower-tail test is appropriate for compactness (lower = more compact
    = more extreme). For partisan efficiency gap (which can be extreme in either
    direction), use the two-sided version: $p = 2 \times \min(\hat{p}, 1-\hat{p})$
    with ESS correction.
\end{enumerate}

The permutation test framework completes the G-track's statistical foundation:
G.7 provides the ensemble (SMC for calibrated inference), G.4 provides the ESS
estimate, and S.1 provides the hypothesis test that converts the ensemble into
a legally interpretable p-value.
